% R911: gerado; alterar dados e gerador, não esta prancha.
\clearpage
\refstepcounter{figure}\label{fig:AR05-flow}
\noindent{\large\sffamily\bfseries\color{navy}AR05 · Fluxograma}\par\smallskip
\noindent\textbf{Ciclo de vida de agentes e perfis gerenciados}\par
\noindent\textit{Como criar, validar, registrar e preparar um agente sem simular trabalho?}\par\medskip
\begin{center}

\tikzset{
 atlasnode/.style={draw=steel,fill=blueBg,rounded corners=2mm,align=center,
  text width=4cm,inner sep=5pt,font=\fontsize{9.5}{11.8}\selectfont,text=textmain},
 atlasroot/.style={atlasnode,fill=navy,text=white,text width=2.8cm,font=\fontsize{10}{12}\selectfont\bfseries},
 atlasdecision/.style={diamond,aspect=2.1,draw=amber,fill=amberBg,align=center,
  text width=3cm,inner sep=2pt,font=\fontsize{9.5}{11.5}\selectfont},
 atlascall/.style={-Stealth,draw=steel,thick},
 atlasevent/.style={-Stealth,draw=violet,dashed,thick},
 atlasdata/.style={-Stealth,draw=teal,densely dotted,thick},
 atlasstructure/.style={draw=slate,thick},
 atlasboundary/.style={-Stealth,draw=amber,dash dot,thick},
 atlasedgelabel/.style={fill=papel,inner sep=1.3pt,font=\fontsize{9}{10}\selectfont}
}

\begin{tikzpicture}
\node[atlasnode,text width=5cm] (s0) {Reutilizar ou descrever um perfil};
\node[atlasnode,text width=5cm,below=6mm of s0] (s1) {Informar capacidades e contrato};
\draw[atlascall] (s0) -- (s1);
\node[atlasnode,text width=5cm,below=6mm of s1] (s2) {Validar spec, teste e permissões};
\draw[atlascall] (s1) -- (s2);
\node[atlasdecision,below=6mm of s2] (decision) {O perfil satisfaz o contrato gerenciado?};
\draw[atlascall] (s2) -- (decision);
\node[atlasnode,fill=amberBg,text width=3.1cm,right=12mm of decision] (refusal) {\textbf{Recusa ou pendência}\\Recusar publicação ou manter legado};
\draw[atlasboundary] (decision.east) -- node[above,font=\small]{não} (refusal.west);
\node[atlasnode,text width=5cm,below=6mm of decision] (s3) {Persistir e registrar a definição};
\draw[atlascall] (decision) -- node[right,font=\small]{sim} (s3);
\node[atlasnode,text width=5cm,below=6mm of s3] (s4) {Preparar tarefa pelo primário};
\draw[atlascall] (s3) -- (s4);
\node[atlasnode,text width=5cm,below=6mm of s4] (s5) {Vincular trabalho e evidência efetivos};
\draw[atlascall] (s4) -- (s5);
\end{tikzpicture}
\end{center}
{\small\textbf{Figura \thefigure.} Fluxograma --- Ciclo de vida de agentes e perfis gerenciados}\par\smallskip
\noindent\textbf{Como funciona e por quê.} Siga a coluna como uma sequência explicativa da operação delimitada. No losango, compare as condições antes de avançar.\par
\noindent\textbf{Justificativa.} Um editor de livro com permissão de escrita recebe escopo e evidências esperadas. Sua criação válida não prova edição realizada: é preciso correlacionar a tarefa atribuída com alterações e verificações do trabalho efetivo.\par
\noindent\textbf{Limite.} O Core não intercepta automaticamente ferramentas externas que criam subagentes. Registro, ativação e cartão available não significam inferência ou conclusão. Alterações de perfis históricos exigem migração explícita, sem usar o baseline para encobrir novas definições. Recuperação: Corrigir contrato sem alterar legado ocultamente\par\smallskip
\fontecore{marceloclaro/orchestrator.py}{MarceloClaroOrchestrator.agent_profile_action}
\fontecore{agents/lifecycle.py}{AgentLifecycleService.create}
